% ============================================
% Section: Communication and Data Sharing
% ============================================
\label{sec:communication}

The communication layer is the heart of the cooperative system: it must move
photo-realistic Gaussian maps between nodes over bandwidth-limited wireless
links, stay consistent under packet loss and intermittent connectivity, and
scale to many nodes without saturating any single link. This section specifies
the transport, message catalogue, connection lifecycle, submap format, and
bandwidth model.

\subsection{Transport}
\label{sec:transport}

The default middleware is \textbf{DDS / ROS 2}, chosen for zero-copy
intra-host transport, configurable QoS, and multicast-based discovery. Where
DDS multicast is blocked (e.g., across subnets), the layer falls back to
\textbf{gRPC} over TCP. Control messages use CBOR encoding; submap payloads use
a custom binary format described in Section~\ref{sec:submap-format} and
Appendix~\ref{app:submap}.

Three QoS profiles are defined:
\begin{itemize}
  \item \emph{Control} (reliable, keep-last 10): HELLO, ELECTION, LEAVE.
  \item \emph{Submap} (best-effort, keep-last 1): a stale submap is useless, so
    we drop rather than queue.
  \item \emph{Trajectory} (reliable, keep-last 1): pose stream must arrive in
    order, but only the latest matters.
\end{itemize}

\subsection{Message Catalogue}
\label{sec:messages}

Table~\ref{tab:messages} lists the message types, their direction, frequency,
payload, and reliability. The catalogue is intentionally small: every message
has a clear owner and a single responsibility, which simplifies correctness
reasoning.

\begin{table}[t]
  \centering
  \caption{Communication message catalogue.}
  \label{tab:messages}
  \small
  \begin{tabular}{lllll}
    \toprule
    Type & Direction & Freq. & Payload & Rel. \\
    \midrule
    HELLO        & node$\to$net   & on join   & NodeDescriptor    & rel. \\
    HEARTBEAT    & node$\to$net   & 5 Hz      & id, pose, health  & BE   \\
    POSE\_STREAM & node$\to$head  & 20 Hz     & spline knots      & rel. \\
    SUBMAP\_PUSH & node$\to$head  & on close  & compressed submap & BE   \\
    SUBMAP\_PULL & head$\to$node  & on demand & bbox query        & rel. \\
    CONSTRAINT   & any$\to$any    & on loop   & SE3 + info        & rel. \\
    MAP\_DELTA   & head$\to$node  & 1 Hz      & merged deltas     & BE   \\
    ELECTION     & node$\to$net   & timeout   & candidate score   & rel. \\
    LEAVE        & node$\to$net   & on exit   & id                & rel. \\
    \bottomrule
  \end{tabular}
\end{table}

\subsection{Connection Lifecycle}
\label{sec:lifecycle}

\paragraph{Join.} A new node broadcasts HELLO with its signed descriptor.
Existing nodes reply with their descriptors; the cluster head replies with the
current cluster map digest and its $\text{node\_id}$.

\paragraph{Calibration handshake.} If the new node's local world frame $W_i$
is unaligned with the head's $W_h$, the head sends an initial prior
$T_{W_h \leftarrow W_i}$ derived from a coarse GNSS / WiFi-RSSI / extrinsic
prior, refined after the first submap exchange (Section~\ref{sec:registration}).

\paragraph{Steady state.} The node streams POSE\_STREAM at 20~Hz and publishes
SUBMAP\_PUSH on submap boundaries---every $S_{\text{sub}}$ meters of travel or
$T_{\text{sub}}$ seconds, whichever comes first (defaults $S_{\text{sub}}=5$~m,
$T_{\text{sub}}=3$~s).

\paragraph{Link loss.} A node that loses the head for more than $3$~s enters
\emph{autonomous mode}: it continues local Gaussian-LIC2, buffers submaps
locally, and re-broadcasts HELLO every $1$~s. On reconnect, buffered submaps
are flushed in chronological order. This satisfies the robustness goal: map
quality stays within 5\% of the connected baseline after a 30~s outage.

\paragraph{Leave / eviction.} LEAVE marks the departing node's submaps as
``owned-but-absent''; they are kept for $T_{\text{keep}}=60$~s then garbage
collected unless referenced by a loop constraint.

\subsection{Submap Format}
\label{sec:submap-format}

A \emph{submap} is the unit of inter-node data sharing: a spatially bounded
subset of a node's Gaussian map, exported on a submap boundary. It consists of
a CBOR header carrying the submap id, owning node id, bounding box in $W_i$,
the body pose at submap close, and a CRC32, followed by a structure-of-arrays
Gaussian blob with means (f16, quantized to 1~cm in the submap local frame),
log-scales, quaternion rotations, logit opacities, spherical harmonics
coefficients (degree-capped), and optional semantic ids. The full binary
layout is given in Appendix~\ref{app:submap}.

\subsection{Compression and Budget Control}
\label{sec:compression}

The compressor exposes four knobs: SH degree (0--3), opacity prune threshold
(0.005--0.05), spatial quantization (1--5~cm), and entropy coding (zstd level
1--9). A \emph{budget controller} monitors a moving average of uplink
throughput and picks the tightest combination that fits the per-node budget
$B_{\text{budget}}$ for the current submap.

\begin{table}[t]
  \centering
  \caption{Compression knobs and their effect.}
  \label{tab:knobs}
  \small
  \begin{tabular}{lll}
    \toprule
    Knob & Range & Effect \\
    \midrule
    SH degree        & 0--3      & halves bytes/Gaussian per degree step \\
    Opacity prune    & 0.005--0.05 & drops 1--5\% of Gaussians, $\approx$lossless \\
    Spatial quant.   & 1--5 cm   & mean quantization error budget \\
    Entropy coding   & zstd 1--9 & 2--4$\times$ on top of above \\
    \bottomrule
  \end{tabular}
\end{table}

On a 20~Mbps Wi-Fi~6 link, a 5~m-travel submap with $\sim$50k Gaussians
compresses to $\approx$3~MB at SH degree 2, 2~cm quantization, zstd-5, i.e.
$\approx$1.2~s at budget---consistent with the real-time target.

\subsection{Bandwidth Model}
\label{sec:bandwidth}

Let $N$ be the number of nodes, $g$ the per-submap Gaussian count, and $f_s$
the submap publication rate. Per-node uplink is
\begin{equation}
  B_{\text{up}} = f_s \cdot \text{sizeof}(\text{compressed\_submap}(g)) \approx f_s \cdot g \cdot 24~\text{B},
\end{equation}
where 24~B is the compressed per-Gaussian cost at SH degree 2. A cluster head's
downlink load is $(N-1) \cdot B_{\text{up}}$. For $N=8$, $g=50$k, $f_s=0.33$~Hz,
the head sees $\approx$3.2~MB/s $\approx$26~Mbps. The head therefore requires a
Wi-Fi~6 or 5G link, while member uplinks fit a 20~Mbps channel. Crucially,
because the number of cluster heads grows as $\mathcal{O}(\sqrt{N})$ rather
than $\mathcal{O}(N)$, total system bandwidth grows sub-linearly in $N$, which
satisfies the scalability goal.

\subsection{Communication Sequence}
\label{sec:sequence}

Figure~\ref{fig:sequence} illustrates the message exchange for a typical
join--publish--merge--loop cycle.

\begin{figure}[t]
  \centering
  \resizebox{\columnwidth}{!}{% Communication sequence diagram for Co-LIC2 paper
% Join -> publish -> merge -> loop -> PGO correction
\begin{tikzpicture}[
  font=\scriptsize,
  >=Stealth,
  actor/.style={draw, rounded corners, minimum width=2.2cm, minimum height=0.55cm, align=center, fill=gray!10},
  msg/.style={->, thick},
  msgd/.style={->, thick, dashed},
]

% Actors (vertical lifelines)
\node[actor] (n) at (0,0) {New Node};
\node[actor] (h) at (4.5,0) {Cluster Head};
\node[actor] (o) at (9,0) {Other Nodes};

\draw[dashed, gray] (n.south) -- ++(0,-6.5);
\draw[dashed, gray] (h.south) -- ++(0,-6.5);
\draw[dashed, gray] (o.south) -- ++(0,-6.5);

% Time flows downward; messages
% Join
\draw[msg] (0,-0.7) -- (4.5,-0.9) node[midway, above, font=\scriptsize] {1. HELLO (NodeDescriptor)};
\draw[msg] (4.5,-1.3) -- (0,-1.5) node[midway, above, font=\scriptsize] {2. cluster digest + head id};

% Publish
\draw[msg] (0,-2.2) -- (4.5,-2.4) node[midway, above, font=\scriptsize] {3. POSE\_STREAM @20Hz};
\draw[msg] (0,-2.9) -- (4.5,-3.1) node[midway, above, font=\scriptsize] {4. SUBMAP\_PUSH (compressed)};

% Merge + delta
\draw[msg] (4.5,-3.7) -- (0,-3.9) node[midway, above, font=\scriptsize] {5. MAP\_DELTA (merged Gaussians)};
\draw[msg] (4.5,-4.3) -- (9,-4.5) node[midway, above, font=\scriptsize] {5'. MAP\_DELTA to peers};

% Loop closure
\draw[msgd, green!50!black] (9,-5.0) -- (4.5,-5.2) node[midway, above, font=\scriptsize] {6. CONSTRAINT (visual loop)};

% PGO correction
\draw[msg, purple!60!black] (4.5,-5.8) -- (0,-6.0) node[midway, above, font=\scriptsize] {7. PGO correction (pose deltas)};
\draw[msg, purple!60!black] (4.5,-6.2) -- (9,-6.4) node[midway, above, font=\scriptsize] {7'. PGO correction};

\end{tikzpicture}
}
  \caption{Communication sequence: a new node joins, publishes submaps, the
  head merges and emits map deltas, a loop closure generates a constraint, and
  the pose-graph optimization broadcasts corrections.}
  \label{fig:sequence}
\end{figure}
